\section{Capítulo~\ref{sec:pain}. El dolor}

Los sensores del dolor, los dos cables, la compuerta en la médula espinal, la perilla descendente de volumen, la sensibilización y la captura de la atención están medidos directamente; las fórmulas de la compuerta y de la sensibilización son modelos simplificados del libro; la regla con la que la máquina elige el volumen de la alarma es una hipótesis.

{\footnotesize
\setlength{\tabcolsep}{3pt}\renewcommand{\arraystretch}{1.15}
\begin{longtable}{>{\raggedright}p{0.5cm}>{\raggedright}p{4.0cm}>{\raggedright}p{5.6cm}>{\raggedright}p{2.3cm}>{\raggedright\arraybackslash}p{1.7cm}}
\toprule
N.º & Proposición & Experimento y qué se midió & Fuente & Estado \\
\midrule\endfirsthead
\toprule
N.º & Proposición & Experimento y qué se midió & Fuente & Estado \\
\midrule\endhead
1 & Umbral alto: los umbrales de calor de los distintos sensores del dolor van de $41^\circ$ a $46$--$48^\circ$ & Registro de fibras individuales. Piel de mono: umbral mediano de los sensores lentos (C), $41^\circ$; de los sensores rápidos de tipo~II, $46^\circ$; de tipo~I, más de $53^\circ$. Piel humana: sensores C que responden también a la presión, $40.7^\circ$; los que no responden a la presión, $48^\circ$. & \cite{Treede1995heat, Weidner1999cfib} & medido \\
2 & Los sensores del dolor se adaptan mucho menos que los del tacto y tras una lesión leve se vuelven más sensibles; el debilitamiento tras un calentamiento intenso se midió en un tipo & Quemadura leve de la piel ($50^\circ$, $100\,\text{s}$): a los $5$--$10$~min el umbral del dolor en personas y el umbral de los sensores C en el mono bajan $2$--$6^\circ$; otros sensores de la piel no muestran ese desplazamiento. En los sensores rápidos de tipo~II la respuesta tras un calentamiento intenso queda suprimida. La comparación de la adaptación con los sensores del tacto es una estimación general, sin un experimento propio que la respalde aquí. & \cite{LaMotte1982hyper, Treede1995heat} & medido \\
3 & Dos cables: rápido, $5$--$30$~m/s; lento, $0.5$--$2$~m/s; tiempo de llegada $t=L/v$~\eqref{eq:paintime} & Velocidad de conducción: sensores rápidos del dolor del mono, en promedio $15$ y $25$~m/s (dos tipos); sensores C humanos, $0.8$--$1.0$~m/s. Con $L=1$~m son decenas de milisegundos y alrededor de un segundo. & \cite{Treede1995heat, Weidner1999cfib} & medido \\
4 & El dolor llega dos veces: primero rápido y preciso, luego sordo y prolongado & Tiempo de reacción al calentamiento del antebrazo humano: de $1100$ a $400\ms$ al subir la temperatura; un calentamiento intenso de la piel con vello produce doble dolor, y el de la palma, no. El primer dolor solo pueden llevarlo fibras de más de $6$~m/s; por su latencia corresponde a los sensores rápidos de tipo~II, que la palma no tiene. El reparto de papeles (“dónde” y “qué tan grave”) es una interpretación. & \cite{CampbellLaMotte1983first, Treede1995heat} & medido \\
5 & Compuerta: la neurona de salida de la médula espinal recibe dolor y tacto, y el tacto excita a la intermediaria inhibidora (fig.~\ref{fig:gate}) & El esquema de la compuerta se propuso como teoría. Ratones, marcaje genético y eliminación de grupos de neuronas: las neuronas excitadoras con somatostatina son necesarias para el dolor mecánico; las neuronas inhibidoras con dinorfina impiden que la entrada de los sensores del tacto llegue a ellas, y sin estas neuronas un roce leve provoca dolor. & \cite{MelzackWall1965, Duan2014} & medido \\
6 & El tacto atenúa el dolor & Personas, pulsos láser dolorosos en la mano: un tacto simultáneo reduce la valoración del dolor y la capacidad de distinguir la energía de los pulsos; el efecto se debilita con la distancia entre el punto del tacto y el del dolor dentro de un mismo segmento cutáneo. & \cite{Mancini2014touch} & medido \\
7 & Supuesto de la teoría de la compuerta: el dolor fuerte apaga por sí mismo a la intermediaria inhibidora, y la compuerta se abre de par en par & En el capítulo se señala como supuesto de la teoría original. El trabajo en ratones describe cómo la compuerta impide que el tacto entre en el canal del dolor; no muestra que el dolor apague a la intermediaria. & \cite{MelzackWall1965} & hipótesis \\
8 & Compuerta~\eqref{eq:gate}: umbral de $0.18$ sin tacto, $0.86$ con tacto, $0.11$ con $g=2$; con $\nu=1$ el tacto reduce la salida de $1$ a $0.25$, con $\nu=2$ no actúa; el tacto solo no pasa (fig.~\ref{fig:gatecurves}) & El cálculo con \texttt{models/\allowbreak pain\_\allowbreak gate.py} con los pesos del texto reproduce todos estos números. & \texttt{models/\allowbreak pain\_\allowbreak gate.py} & modelo \\
9 & Las vías descendentes del tronco cambian la ganancia de la compuerta en ambos sentidos & Rata, bulbo raquídeo, registro durante el calentamiento de la cola: unas neuronas descargan antes de la retirada (“on”), otras se callan (“off”); una estimulación débil de esta zona suprime el reflejo. Revisión: los mismos circuitos facilitan e inhiben la transmisión del dolor, y los opioides actúan a través de ellos. & \cite{Fields1983rvm, Fields2004} & medido \\
10 & La expectativa de alivio atenúa el dolor a través del canal opioide & Personas con dolor agudo: en aquellas a quienes la expectativa de alivio les redujo el dolor, el bloqueo de los receptores opioides devolvió el dolor al nivel de las demás; en las demás, el bloqueo no cambió el dolor. & \cite{Levine1978} & medido \\
11 & El cerebro elige el volumen de la alarma según el beneficio de la interrupción & No hay experimento que haya medido la regla de elección del volumen. & --- & hipótesis \\
12 & Sensibilización: tras una lesión, la salida de la compuerta crece y el umbral baja, en parte por la propia médula espinal; persiste después de que cesa el estímulo lesivo & Rata: tras una lesión de la piel, el umbral del reflejo flexor baja y la respuesta crece; el análisis electrofisiológico muestra que parte del aumento surge en la propia médula espinal. El estímulo lesivo provoca alteraciones duraderas de la sensibilidad que sobreviven al propio estímulo. El capítulo no da el tiempo exacto de recuperación. & \cite{Woolf1983} & medido \\
13 & La sensibilización~\eqref{eq:sensitization} es una realimentación positiva; con un lazo fuerte, la alarma puede trabarse después de la curación & Se deduce del modelo~\eqref{eq:sensitization} (ejercicio al final del capítulo). No hay experimento que muestre que el dolor persistente tras la curación sea un lazo trabado precisamente de este tipo. & deducción en el capítulo & hipótesis \\
14 & El dolor es una recompensa negativa, y el aprendizaje con él sigue el error de diferencias temporales & Resonancia, personas, cadenas de precursores del dolor: la actividad del estriado ventral y de la ínsula anterior coincide con las señales de aprendizaje por diferencias temporales. Mono: parte de las neuronas \nt{DOPA} se excita con un soplo de aire desagradable y su precursor; otra parte se inhibe. & \cite{Seymour2004, Matsumoto2009} & medido \\
15 & El dolor interrumpe el trabajo en curso & Personas, estímulo eléctrico doloroso y sondas sonoras: la respuesta a una sonda $100\ms$ después del inicio del dolor es notablemente más lenta; a los $1.5\,\text{s}$ ya no hay diferencia con el control. Una revisión reúne estos experimentos en un modelo de captura de la atención. & \cite{Crombez1996disrupt, EcclestonCrombez1999} & medido \\
16 & La retirada se cierra en la médula espinal & Rata: las neuronas de las capas profundas del asta dorsal reproducen el patrón espacial del reflejo de retirada de cada músculo; no se logra excitarlas antidrómicamente desde el segmento cervical, es decir, son intermediarias espinales. & \cite{Schouenborg1995wr} & medido \\
\bottomrule
\end{longtable}}
